% =================================================================
%  TechnicalReport.tex
%  Technical Report: KPI Framework for the 2026 FIFA World Cup
%                    Group-Stage Schedule Evaluation
%  CORS OR Challenge 2026 -- Team UofT Champions
%
%  Revision notes (this version):
%   - Travel KPIs (1.1-1.4) reformulated for base-camp logistics
%   - KPI 1.7 (Entry and Visa Restriction Exposure) added
%   - KPI 2.4 (Altitude) reformulated for base-camp model
%   - KPI 3.3: positional advantage score a_i defined exactly
%   - Summary table updated (25 KPIs total)
% =================================================================

\documentclass[12pt, letterpaper]{article}

% --- Page geometry -----------------------------------------------
\usepackage[top=1in, bottom=1in, left=1in, right=1in]{geometry}

% --- Math --------------------------------------------------------
\usepackage{amsmath}
\usepackage{amssymb}
\usepackage{mathtools}

% --- Typography --------------------------------------------------
\usepackage[T1]{fontenc}
\usepackage{lmodern}
\usepackage{microtype}
\usepackage{setspace}

% --- Coloured boxes for KPIs ------------------------------------
\usepackage[most]{tcolorbox}

% --- Tables ------------------------------------------------------
\usepackage{booktabs}
\usepackage{longtable}
\usepackage{array}

% --- Headers and footers ----------------------------------------
\usepackage{fancyhdr}

% --- Lists -------------------------------------------------------
\usepackage{enumitem}

% --- References (author--year, parenthetical) -------------------
\usepackage[round, authoryear, sort&compress]{natbib}

% --- Hyperlinks --------------------------------------------------
\usepackage[hidelinks, breaklinks=true]{hyperref}

% =================================================================
%  Custom KPI box
%  Usage: \begin{kpibox}{KPI X.Y --- Title}  ...  \end{kpibox}
% =================================================================
\newtcolorbox{kpibox}[1]{
  enhanced,
  breakable,
  colback      = gray!5!white,
  colframe     = black!40,
  fonttitle    = \bfseries\small,
  coltitle     = black,
  colbacktitle = gray!18!white,
  title        = {#1},
  boxrule      = 0.65pt,
  left  = 6pt,
  right = 6pt,
  top   = 4pt,
  bottom= 4pt,
  toptitle    = 3pt,
  bottomtitle = 3pt,
}

% Short-cut for field labels inside KPI boxes
\newcommand{\kpilabel}[1]{\noindent\textbf{#1}\enspace}

% =================================================================
%  Headers and footers
% =================================================================
\pagestyle{fancy}
\fancyhf{}
\lhead{\small KPI Framework \textendash{} 2026 FIFA World Cup Group-Stage Schedule}
\rhead{\small CORS OR Challenge 2026}
\rfoot{\small \thepage}
\renewcommand{\headrulewidth}{0.4pt}

% =================================================================
%  Spacing
% =================================================================
\onehalfspacing
\setlength{\parskip}{3pt}

% =================================================================
%  Title
% =================================================================
\title{%
  \LARGE\textbf{CORS OR Challenge 2026 \textendash{} Technical Report}
}
\author{%
  \textbf{Team UofT Champions}\\[4pt]
  \small University of Toronto
}
\date{May 2026}

% =================================================================
\begin{document}
% =================================================================

\maketitle
\thispagestyle{empty}
\vspace{10pt}

\vspace{6pt}
%\tableofcontents
\newpage

% =================================================================
%  1.  INTRODUCTION
% =================================================================
\section{Key Performance Indicators for Evaluating the 2026 FIFA
         World Cup Group-Stage Schedule}

This report defines the key performance indicators (KPIs) used
to evaluate the group-stage schedule of the 2026 FIFA World Cup as part
of the CORS OR Challenge 2026.  The group stage comprises 72 matches
distributed across 16 host cities in Canada, Mexico, and the United States
between June~11 and June~27, 2026.  We organise our evaluation framework
around \textbf{five categories of scheduling criteria}: (1)~\emph{Travel,
Relocation, and Rest}, (2)~\emph{Heat and Environmental Exposure},
(3)~\emph{Competitive Integrity}, (4)~\emph{Operational Efficiency}, and
(5)~\emph{Broadcast Attractiveness and Commercial Value}.  Within these five
categories, we propose 19 individual KPIs together with 5 composite
cross-criteria KPIs and 1 robustness indicator, for a total of 25 formally
defined performance measures.  Each KPI is accompanied by a formal
description, a precise mathematical formula, and a citation.

\subsection{Five Criteria Categories}

\begin{enumerate}[leftmargin=*, label=\textbf{C\arabic*.}, itemsep=6pt]

  \item \textbf{Travel, Relocation, and Rest.}
        How far teams must travel on round trips between their base camp
        and each group-stage venue; whether the travel burden is equitably
        shared within each group; the jet-lag cost of crossing time zones;
        the geographic spread of a team's match venues across clusters;
        whether all teams receive adequate and symmetric rest between
        consecutive matches; and the impact of US travel-ban and visa
        policies on fan accessibility at US-hosted fixtures.

  \item \textbf{Heat and Environmental Exposure.}
        The degree to which the schedule exposes players to dangerous heat
        (quantified via the Wet-Bulb Globe Temperature, WBGT); whether
        heat-risk matches are assigned to roofed/climate-controlled venues;
        and the performance disruption caused by large elevation changes
        between a team's base camp and high-altitude Mexican venues.

  \item \textbf{Competitive Integrity.}
        Whether the schedule permits or discourages collusion in final-round
        group matches; the probability and severity of stakeless
        (``dead-rubber'') matches; and whether the ordering of matches within
        a group systematically advantages particular teams through first-mover
        effects.

  \item \textbf{Operational Efficiency.}
        The evenness of the match-and-operations load across the 16 host
        venues, and the degree to which the schedule enables fans to attend
        affordably.

  \item \textbf{Broadcast Attractiveness and Commercial Value.}
        Alignment of kickoff times with peak viewing windows in key global
        markets; the quality of time slots allocated to high-profile matches;
        and the equitable distribution of commercially valuable fixtures
        across host cities.

\end{enumerate}

\subsection{Importance Ranking and KPI Density}
\label{sec:ranking}

The five criteria carry different analytical weights for the 2026 edition.
The KPI density in this report reflects the following judgment.

\textbf{C1 (Travel/Rest) and C2 (Heat)} are the most consequential
criteria.  For heat, \citet{MullanEtAl2025} find that 14 of 16 host cities
frequently exceed 28\,$^\circ$C WBGT in afternoon slots, and FIFPRO identifies
six venues as ``extremely high risk''~\citep{FIFPRO2025}.  For travel,
2026's continent-spanning geography is a direct instantiation of the
Traveling Tournament Problem~\citep{Easton2001}, with documented intra-group
travel disparities exceeding 2{,}000~km~\citep{MikamiTravel2026}.  Both
criteria receive four to seven KPIs each.

\textbf{C3 (Competitive Integrity)} carries the deepest academic literature
but the 12$\times$4 format already resolves the primary collusion concern by
design~\citep{GuajardoKrumer2024, FIFA2023Format}; three targeted KPIs
suffice.

\textbf{C4 (Operational Efficiency) and C5 (Broadcast)} are genuine
scheduling drivers but subordinate to welfare and fairness in a multi-objective
programme.  They receive two to three KPIs each and primarily serve as
tension terms.

In addition to the individual KPIs, Section~\ref{sec:composite} proposes five
\emph{composite KPIs} spanning multiple criteria, including a
Recovery-Adjusted Hardship Index (RAHI) and a Schedule Fairness Index (SFI).

% =================================================================
%  2.  NOTATION
% =================================================================
\section{Mathematical Notation}

The following symbols are used throughout the report.  All distances are
great-circle (Haversine) distances in kilometres; $\mathrm{dist}(\cdot,\cdot)$
applies to any two geographic points, not only venues in $V$.
Rest days are calendar days.

\begin{center}
\renewcommand{\arraystretch}{1.30}
\begin{tabular}{@{}lp{10cm}@{}}
\toprule
\textbf{Symbol} & \textbf{Definition} \\
\midrule
$T$ & Set of 48 teams; generic index $i$ \\
$G$ & Set of 12 groups (each $|g|=4$); generic index $g$ \\
$M$ & Set of 72 group-stage matches; generic index $m$ \\
$V$ & Set of 16 host venues; generic index $v$ \\
$D$ & Set of calendar dates June~11--June~27; generic index $d$ \\
$S$ & Set of available kickoff slots (date$\times$time pairs) \\
$v_{i,k}$ & Venue of team $i$'s $k$-th group match,
            $k\in\{1,2,3\}$ \\
$d_{i,k}$ & Date of team $i$'s $k$-th group match \\
$s_{i,k}$ & Kickoff slot of team $i$'s $k$-th group match \\
$g_{i,k}$ & Rest days between matches $k$ and $k{+}1$ of team $i$:
            $\;g_{i,k} = d_{i,k+1} - d_{i,k} - 1$ \\
$b_{i}$ & Base camp of team $i$ (a fixed geographic location,
          not necessarily in $V$) \\
$\mathrm{dist}(u,v)$ & Great-circle distance between two geographic
                       points $u$ and $v$ (km) \\
$\mathrm{tz}(v)$ & UTC offset of venue $v$ (hours) \\
$\mathrm{tz}(b_i)$ & UTC offset of team $i$'s base camp (hours) \\
$\mathrm{elev}(v)$ & Elevation of venue $v$ (metres above sea level) \\
$\mathrm{elev}(b_i)$ & Elevation of team $i$'s base camp (metres) \\
$\mathrm{cluster}(v)$ & FIFA zone of venue $v$: West, Central, or East \\
$R \subseteq V$ & Roofed or climate-controlled venues \\
$V^{\mathrm{US}} \subseteq V$ & US-hosted venues \\
$r(m,i)$ & Rest days of team $i$ immediately before match $m$ \\
$E_i$ & FIFA/Elo rating of team $i$ (lower $=$ stronger) \\
$P \subseteq M$ & High-profile matches (e.g., both teams ranked
                  FIFA top-20) \\
$F^{\mathrm{ban}} \subseteq T$ & Teams from countries subject to a
                                 full US entry ban \\
$F^{\mathrm{bond}} \subseteq T$ & Teams from countries subject to the
                                  US visa bond programme \\
\bottomrule
\end{tabular}
\end{center}

% =================================================================
%  3.  CRITERION 1: TRAVEL, RELOCATION, AND REST
% =================================================================
\section{Criterion 1: Travel, Relocation, and Rest}

\subsection{Overview}

The geographic scale of 2026 makes travel the single largest logistical
challenge of the group stage.  Venues span four time zones and three
countries, and group-stage fixture assignments can require teams to cover
more than 5{,}000~km of round-trip distance across their three
matches~\citep{SITravel2026}.

\medskip
\textbf{Base-camp logistics.}
A defining feature of 2026 that distinguishes it from prior tournaments
is the adopted \emph{base-camp model}: each team establishes a single
fixed base camp $b_i$ at a FIFA-approved city before the group stage
begins, departs for each match, and returns to $b_i$ immediately
afterwards.  Unlike the classical ``stay at venue'' model---in which
teams travel directly between consecutive venues and rest at each
destination---the base-camp model means that every match generates an
independent round trip $b_i \!\to\! v_{i,k} \!\to\! b_i$.
All travel-related KPIs in this section are formulated under the
base-camp assumption.

\medskip
\textbf{Rest inequality.}
Recovery between matches is a direct welfare concern: insufficient rest
impairs physical performance and increases injury risk, while asymmetric
rest between opponents introduces a structural competitive advantage
unrelated to playing ability.

\medskip
\textbf{Entry and visa restrictions.}
A novel 2026-specific concern: US government travel bans and visa bond
requirements affecting nationals of several qualified nations.

\subsection{KPIs}

% ----- KPI 1.1 -----
\begin{kpibox}{KPI~1.1 \textendash{} Total Per-Team Travel Distance}

\kpilabel{Description.}
The total distance a team travels across its three group-stage matches
under the base-camp model.  Each match $k$ generates a round trip from
the base camp $b_i$ to venue $v_{i,k}$ and back, so there are three
independent round trips (one per match) rather than two inter-venue legs.
This is formulated in the classical Traveling Tournament Problem
formulation~\citep{Easton2001} (which assumed teams stay at each venue
and travel directly to the next).  Summing across all 48 teams gives the
tournament-level efficiency objective; minimising the maximum gives the
equity objective.

\medskip
\kpilabel{Formula.}
\begin{align*}
  \mathrm{TD}_i &= 2\sum_{k=1}^{3} \mathrm{dist}(b_i,\, v_{i,k}),
  \qquad i \in T.
\end{align*}
The factor of~2 accounts for both the outbound and return legs of each
round trip.

\medskip\noindent
Efficiency objective: $\;\displaystyle\min_{S}\;\sum_{i \in T}
\mathrm{TD}_i$.
\quad
Equity (minimax) objective: $\;\displaystyle\min_{S}\;\max_{i \in T}
\mathrm{TD}_i$.

\medskip
\kpilabel{Reference.}
\citet{Easton2001} (classical TTP); \citet{GuajardoEtAl2026};
\citet{SITravel2026}.  Base-camp adaptation is original (this work).
\end{kpibox}

% ----- KPI 1.2 -----
\begin{kpibox}{KPI~1.2 \textendash{} Intra-Group Travel Dispersion}

\kpilabel{Description.}
The inequality of base-camp round-trip travel burdens \emph{within} a
group.  Because group-mates compete directly for the same qualification
slots, a large disparity in total travel constitutes an unearned
competitive advantage.  \citet{GuajardoEtAl2026} formally frame this
as the fairness complement of KPI~1.1, noting disparities such as
Argentina ($\approx700$~km total) vs.\ Algeria ($\approx2{,}400$~km)
within the same group---inequalities that arise from uneven venue
assignments regardless of base-camp choices.

\medskip
\kpilabel{Formula.}
Using the base-camp total travel $\mathrm{TD}_i$ from KPI~1.1,
define for group $g$ the range
\begin{align*}
  \Delta_g &= \max_{i \in g}\,\mathrm{TD}_i \;-\;
              \min_{i \in g}\,\mathrm{TD}_i
\end{align*}
and the coefficient of variation
\begin{align*}
  \mathrm{CV}_g &=
    \frac{\sigma\!\bigl(\{\mathrm{TD}_i : i \in g\}\bigr)}
         {\mu\!\bigl(\{\mathrm{TD}_i : i \in g\}\bigr)}.
\end{align*}
\vspace{-8pt}

\noindent
Fairness objectives: $\;\displaystyle\min_{S}\;\sum_{g \in G}\Delta_g$
\;or\; $\;\displaystyle\min_{S}\;\max_{g \in G}\Delta_g$.

\medskip
\kpilabel{Reference.}
\citet{GuajardoEtAl2026}; \citet{MikamiTravel2026}.
\end{kpibox}

% ----- KPI 1.3 -----
\begin{kpibox}{KPI~1.3 \textendash{} Circadian Shift Cost (Jet-Lag Burden)}

\kpilabel{Description.}
Under the base-camp model, every team returns to its base camp $b_i$
immediately after each match and spends its rest days there.  Because
players are continuously re-exposed to the base-camp light--dark cycle
between matches, their circadian clock remains anchored to
$\mathrm{tz}(b_i)$ throughout the group stage.  The performance-relevant
jet-lag effect is therefore about the mismatch between the \emph{kickoff time at the
match venue} and the players' internal body clock, which is set to
base-camp local time.

Specifically, if the match at venue $v_{i,k}$ kicks off at local time
$\tau_{i,k}$ (e.g., 13:00 local), the players' bodies perceive that
kickoff at the equivalent base-camp clock time
\begin{align*}
  \hat{\tau}_{i,k} &= \tau_{i,k}
    - \bigl[\mathrm{tz}(v_{i,k}) - \mathrm{tz}(b_i)\bigr],
\end{align*}
expressed modulo 24~hours.  A kickoff that falls during the players'
subjective night window (approximately 23:00--07:00 base-camp time)
constitutes a meaningful circadian penalty: alertness, reaction time,
and peak muscular output are significantly reduced during those hours
even when players appear awake~\citep{Recht1995}.

\medskip
\kpilabel{Formula.}
Let $\tau_{i,k} \in [0,24)$ denote the kickoff time of match $k$ of
team $i$, expressed in local venue time (hours).  Define the
\emph{perceived kickoff time} in base-camp time:
\begin{align*}
  \hat{\tau}_{i,k} &= \Bigl(\tau_{i,k}
    - \bigl[\mathrm{tz}(v_{i,k}) - \mathrm{tz}(b_i)\bigr]\Bigr)
    \bmod 24.
\end{align*}
Define the \emph{circadian penalty function} measuring how far the
perceived kickoff falls inside the subjective night window
$[\tau_{\mathrm{lo}},\, \tau_{\mathrm{hi}}]$
(default: $\tau_{\mathrm{lo}} = 23$, $\tau_{\mathrm{hi}} = 7$,
i.e., the window wraps around midnight):
\begin{align*}
  \phi(\hat{\tau}) &= \begin{cases}
    \min\!\bigl(\hat{\tau} - \tau_{\mathrm{lo}},\;
                \tau_{\mathrm{hi}} + 24 - \tau_{\mathrm{lo}}\bigr)
      & \text{if } \hat{\tau} \geq \tau_{\mathrm{lo}}, \\[4pt]
    \min\!\bigl(\hat{\tau} + 24 - \tau_{\mathrm{lo}},\;
                \tau_{\mathrm{hi}} - \hat{\tau}\bigr)
      & \text{if } \hat{\tau} \leq \tau_{\mathrm{hi}}, \\[4pt]
    0 & \text{otherwise,}
  \end{cases}
\end{align*}
so $\phi(\hat{\tau}) \in [0,8]$ is the depth of intrusion (in hours)
into the subjective night, and $\phi = 0$ whenever the perceived
kickoff falls entirely within the subjective day.

The cumulative circadian disruption score for team $i$ across its three
group-stage matches is:
\begin{align*}
  \mathrm{JL}_i &= \sum_{k=1}^{3} \phi\!\bigl(\hat{\tau}_{i,k}\bigr).
\end{align*}
\vspace{-8pt}

\noindent
Objective: $\;\displaystyle\min_{S}\;\sum_{i \in T}\mathrm{JL}_i$.
\quad
Fairness objective: $\;\displaystyle\min_{S}\;\max_{i \in T}\mathrm{JL}_i$.

\medskip
\kpilabel{Reference.}
\citet{Recht1995} (circadian performance decrement of 2--7\%;
eastward harder than westward); \citet{KendallEtAl2010}.
Perceived-kickoff formulation is original (this work).
\end{kpibox}

% ----- KPI 1.4 -----
\begin{kpibox}{KPI~1.4 \textendash{} Match-Venue Geographic Dispersion}

\kpilabel{Description.}
How many distinct FIFA geographic clusters
(West/Central/East) must a team visit across its three match venues?
A team whose three venues all fall within the same cluster can select a
centrally located base camp and execute short round trips, while a team
whose venues span all three clusters faces at least one long-haul
transcontinental flight regardless of base-camp placement.  A
2026-specific variant counts the number of individual match trips that
cross an international border from the base camp, triggering additional
visa, customs, and security obligations unique to this three-nation
tournament.

\medskip
\kpilabel{Formula.}
Venue cluster diversity:
\begin{align*}
  \mathrm{CC}_i &= \bigl|\bigl\{\mathrm{cluster}(v_{i,1}),\,
    \mathrm{cluster}(v_{i,2}),\,
    \mathrm{cluster}(v_{i,3})\bigr\}\bigr|
  \;\in\;\{1,2,3\}.
\end{align*}
Base-camp border-crossing count:
\begin{align*}
  \mathrm{BC}_i &= \bigl|\bigl\{k \in \{1,2,3\} :
    \mathrm{country}(v_{i,k}) \neq
    \mathrm{country}(b_i)\bigr\}\bigr|
  \;\in\;\{0,1,2,3\}.
\end{align*}
\vspace{-8pt}

\noindent
Primary objective: $\;\displaystyle\min_{S}\;\sum_{i \in T}\mathrm{CC}_i$.

\medskip
\kpilabel{Reference.}
\citet{MikamiTravel2026}; \citet{KendallEtAl2010}.
Reformulation for base-camp logistics is original (this work).
\end{kpibox}

% ----- KPI 1.5 -----
\begin{kpibox}{KPI~1.5 \textendash{} Per-Team Rest Adequacy}

\kpilabel{Description.}
The minimum number of full calendar days separating any two consecutive
matches of a given team.  FIFA's published schedule design targets at
least three rest days for every gap, consistent with the 72-hour
recovery minimum recommended in elite sports medicine.  This KPI
measures both individual compliance and the margin above the hard
welfare floor.  The rest calculation is independent of the base-camp
model: rest days are determined by match dates alone.

\medskip
\kpilabel{Formula.}
\begin{align*}
  \mathrm{RA}_i &= \min\!\bigl(g_{i,1},\; g_{i,2}\bigr),
  \qquad g_{i,k} = d_{i,k+1} - d_{i,k} - 1.
\end{align*}
\vspace{-8pt}

\noindent
Hard constraint: $g_{i,k} \geq g_{\min}$ for all $i \in T$,
$k \in \{1,2\}$, where $g_{\min} = 2$ (i.e., 72~hours).
Report the deficiency count $N_{\mathrm{short}} =
|\{(i,k) : g_{i,k} < 3\}|$.

\medskip
\kpilabel{Reference.}
FIFA schedule announcement (June~2024): ``three days of rest for
teams is observed for 103 of the tournament's 104 matches.''
\citet{AtanCavdaroglu2018}; \citet{KendallEtAl2010}.
\end{kpibox}

% ----- KPI 1.6 -----
\begin{kpibox}{KPI~1.6 \textendash{} Rest Asymmetry Between Opponents}

\kpilabel{Description.}
Even when every team receives adequate individual rest (KPI~1.5), one
team may enter a particular match meaningfully fresher than its opponent,
constituting a structural competitive advantage unrelated to team quality.
\citet{AtanCavdaroglu2018} define a \emph{rest mismatch} as any match in
which the two opponents had different numbers of rest days, and show that
minimising the total count of such mismatches is an integer-programmable
objective.

\medskip
\kpilabel{Formula.}
For match $m$ between teams $i$ and $j$:
\begin{align*}
  \mathrm{RD}_m &= \bigl|r(m,i) - r(m,j)\bigr|.
\end{align*}
Two objectives:
\begin{align*}
  \text{Count objective:}\quad
    &\min_{S}\;\bigl|\{m \in M : \mathrm{RD}_m > 0\}\bigr|,
  \\[4pt]
  \text{Magnitude objective:}\quad
    &\min_{S}\;\sum_{m \in M} \mathrm{RD}_m.
\end{align*}

\medskip
\kpilabel{Reference.}
\citet{AtanCavdaroglu2018}; \citet{CavdarogluAtan2022}.
\end{kpibox}

% ----- KPI 1.7 -----
\begin{kpibox}{KPI~1.7 \textendash{} Entry and Visa Restriction Exposure (ERI)}

\kpilabel{Description.}
The 2026 World Cup's three-nation hosting structure introduces a
scheduling fairness dimension with no precedent in prior tournaments.
The US government's expanded travel ban (effective January~2026) bars
nationals of several qualified nations from entering the United States.
While players and team staff are exempted under Presidential Proclamation
10998, ordinary fans face the full restriction, making it impossible for
them to attend US-hosted matches~\citep{ESPN2026TravelBan,
NPR2026TravelBan}.  As of May~2026, four qualified nations---Iran,
Haiti, Senegal, and C\^{o}te d'Ivoire---are subject to a full entry
suspension ($F^{\mathrm{ban}}$).  A further three nations---Algeria,
Cabo Verde, and Tunisia---are subject to a US visa bond programme
requiring applicants to post deposits of \$5{,}000--\$15{,}000
($F^{\mathrm{bond}}$), a severe practical barrier for most ordinary
supporters~\citep{AIC2026TravelBan}.

The schedule determines which matches are hosted in the United States
(78~of the 104 group-stage and knockout fixtures), directly controlling
the exposure of affected nations to these restrictions.  Minimising this
KPI is equivalent to concentrating the fixtures of affected nations at
Canadian and Mexican venues wherever flexibility exists.

\medskip
\kpilabel{Formula.}
Assign a severity weight to each team:
\begin{align*}
  \sigma(t) &= \begin{cases}
    1.0 & \text{if } t \in F^{\mathrm{ban}}, \\
    0.5 & \text{if } t \in F^{\mathrm{bond}}, \\
    0   & \text{otherwise.}
  \end{cases}
\end{align*}
For match $m$ between teams $i(m)$ and $j(m)$ at venue $v(m)$, the
entry and visa restriction exposure index is:
\begin{align*}
  \mathrm{ERI} &= \sum_{m \in M}
    \mathbf{1}\!\bigl[v(m) \in V^{\mathrm{US}}\bigr]\cdot
    \max\!\bigl(\sigma(i(m)),\;\sigma(j(m))\bigr).
\end{align*}
Unweighted count of affected US-hosted matches:
\begin{align*}
  N_{\mathrm{ban}} &= \sum_{m \in M}
    \mathbf{1}\!\bigl[v(m) \in V^{\mathrm{US}}\bigr]\cdot
    \mathbf{1}\!\bigl[i(m) \in F^{\mathrm{ban}} \cup F^{\mathrm{bond}}
    \;\vee\; j(m) \in F^{\mathrm{ban}} \cup F^{\mathrm{bond}}\bigr].
\end{align*}
\vspace{-8pt}

\noindent
Objective: $\;\displaystyle\min_{S}\;\mathrm{ERI}$.

\medskip
\kpilabel{Reference.}
\citet{ESPN2026TravelBan}; \citet{NPR2026TravelBan};
\citet{AIC2026TravelBan}; \citet{FIFA2023Format}.
ERI formulation is original (this work).
\end{kpibox}

% =================================================================
%  4.  CRITERION 2: HEAT AND ENVIRONMENTAL EXPOSURE
% =================================================================
\section{Criterion 2: Heat and Environmental Exposure}

\subsection{Background: Wet-Bulb Globe Temperature (WBGT)}
\label{sec:wbgt}

\textbf{Definition and physical interpretation.}
The Wet-Bulb Globe Temperature (WBGT) is the internationally recognised
heat-stress index in both sports medicine and occupational health.
Originally developed by the U.S.\ Army in the 1950s~\citep{YaglouMinard1957},
it is now mandated by the International Olympic Committee, the
International Labour Organization, and the American College of Sports
Medicine.

WBGT outperforms dry-bulb (``shade'') temperature as a welfare metric
because it simultaneously accounts for three physical pathways of heat
stress:
\begin{itemize}[noitemsep, topsep=4pt]
  \item \emph{Humidity}, via the natural wet-bulb temperature
        $T_\mathrm{nwb}$ (weight~0.70 outdoors): the dominant factor,
        because high humidity impedes evaporative cooling through
        perspiration.
  \item \emph{Radiant heat from solar radiation}, via the black-globe
        temperature $T_g$ (weight~0.20 outdoors): relevant only outdoors.
  \item \emph{Air temperature}, via the dry-bulb temperature
        $T_\mathrm{db}$ (weight~0.10 outdoors).
\end{itemize}

For \emph{outdoor, unshaded} conditions:
\begin{equation}
  \mathrm{WBGT}_m
    \;=\; 0.70\,T_\mathrm{nwb} \;+\; 0.20\,T_g \;+\; 0.10\,T_\mathrm{db}.
  \label{eq:wbgt_outdoor}
\end{equation}

For \emph{indoor or fully shaded} conditions (roofed stadiums, no direct
solar radiation):
\begin{equation}
  \mathrm{WBGT}_m
    \;=\; 0.70\,T_\mathrm{nwb} \;+\; 0.30\,T_\mathrm{db}.
  \label{eq:wbgt_indoor}
\end{equation}

When only standard meteorological station data are available
(temperature, relative humidity, dewpoint, solar irradiance),
$T_\mathrm{nwb}$ and $T_g$ may be estimated using the empirical models
of \citet{Liljegren2008}.

\medskip
\textbf{Key action thresholds.}
FIFPRO guidelines~\citep{FIFPRO2025} define:

\begin{center}
\begin{tabular}{@{}lp{8cm}@{}}
\toprule
WBGT & Required action \\
\midrule
$\geq 26\,^\circ$C & Mandatory cooling breaks every 15~minutes \\
$\geq 28\,^\circ$C & Match postponement recommended \\
$\geq 32\,^\circ$C & Extreme risk; immediate action required \\
\bottomrule
\end{tabular}
\end{center}

\medskip
\textbf{2026 context.}
\citet{MullanEtAl2025} analyse ERA5 reanalysis climatology (2003--2022)
for all 16 host cities and find that 14 of 16 cities frequently exceed
$28\,^\circ$C WBGT in afternoon slots.  Six venues---Atlanta, Dallas,
Houston, Kansas City, Miami (USA), and Monterrey (Mexico)---are rated
``extremely high risk'' by FIFPRO~\citep{FIFPRO2025}.  Only five of the
16 stadiums (Atlanta, Dallas, Houston, Los Angeles, Vancouver) possess
a retractable roof or climate-control system, meaning structural
mitigation is unavailable at most
venues~\citep{WorldWeatherAttribution2026}.

\subsection{Background: Altitude and Acclimatisation}

Two Mexican venues present a distinct physiological challenge.
Mexico City (Estadio Azteca) sits at approximately 2{,}240~m above sea
level; Guadalajara (Estadio Akron) at approximately 1{,}560~m.  At such
altitudes, the partial pressure of oxygen is reduced.  Under the base-camp model, a team playing at altitude must
complete a round trip $b_i \!\to\! v_{i,k} \!\to\! b_i$ for each
high-altitude match, so the acclimatisation demand is determined by the
elevation differential between $b_i$ and $v_{i,k}$, not between
consecutive venues.  FIFA's sports-medicine guidelines recommend a
minimum of ten days at elevations above 3{,}000~m; even moderate
elevations of 1{,}500$+$~m can reduce aerobic performance within a
standard three-day rest window.

\subsection{KPIs}

% ----- KPI 2.1 -----
\begin{kpibox}{KPI~2.1 \textendash{} Tournament WBGT Heat-Stress Score}

\kpilabel{Description.}
The aggregate heat-stress load across all 72 group-stage matches.
For each match we compute the expected WBGT at the scheduled venue,
date, and kickoff time using
equations~\eqref{eq:wbgt_outdoor}--\eqref{eq:wbgt_indoor} and site
climatology~\citep{Liljegren2008}.  The KPI also reports exceedance
counts at the three FIFPRO action thresholds.

\medskip
\kpilabel{Formula.}
\begin{align*}
  H_{\mathrm{total}} &= \sum_{m \in M} \mathrm{WBGT}_m.
\end{align*}
Threshold exceedance counts (to be reported separately):
\begin{align*}
  N_{26} &= \bigl|\{m \in M : \mathrm{WBGT}_m \geq
            26\,^\circ\text{C}\}\bigr|,\\
  N_{28} &= \bigl|\{m \in M : \mathrm{WBGT}_m \geq
            28\,^\circ\text{C}\}\bigr|,\\
  N_{32} &= \bigl|\{m \in M : \mathrm{WBGT}_m \geq
            32\,^\circ\text{C}\}\bigr|.
\end{align*}
\vspace{-8pt}

\noindent
Primary objective: $\;\displaystyle\min_{S}\;N_{28}$.

\medskip
\kpilabel{Reference.}
\citet{Liljegren2008} (WBGT formula); \citet{FIFPRO2025} (thresholds);
\citet{MullanEtAl2025} (2026 climatology);
\citet{WorldWeatherAttribution2026}.
\end{kpibox}

% ----- KPI 2.2 -----
\begin{kpibox}{KPI~2.2 \textendash{} Per-Team Heat Load}

\kpilabel{Description.}
The cumulative excess WBGT exposure absorbed by each team across its
three group-stage matches, measured above the 28\,$^\circ$C safety
threshold.  Because opponents are fixed by the draw, the schedule can
still allocate high-risk matches to cooler (evening) slots or covered
venues.  This KPI identifies which teams bear a disproportionately large
heat burden, enabling a fairness assessment analogous to KPI~1.2.

\medskip
\kpilabel{Formula.}
\begin{align*}
  \mathrm{HL}_i &= \sum_{k=1}^{3}
    \max\!\bigl(0,\;\mathrm{WBGT}_{m(i,k)} - \tau\bigr),
  \qquad \tau = 28\,^\circ\text{C}.
\end{align*}
\vspace{-8pt}

\noindent
Fairness (minimax) objective:
$\;\displaystyle\min_{S}\;\max_{i \in T}\,\mathrm{HL}_i$.

\medskip
\kpilabel{Reference.}
\citet{FIFPRO2025}; \citet{MullanEtAl2025}; \citet{Liljegren2008}.
\end{kpibox}

% ----- KPI 2.3 -----
\begin{kpibox}{KPI~2.3 \textendash{} Roof and Climate-Control Conditioning Ratio}

\kpilabel{Description.}
The fraction of safety-threshold-exceeding matches that are assigned to
roofed or climate-controlled venues.  A covered stadium eliminates the
solar radiation term in equation~\eqref{eq:wbgt_outdoor}, switching to
the safer indoor WBGT formula~\eqref{eq:wbgt_indoor} and substantially
reducing player heat exposure.  With only five covered venues hosting a
minority of group fixtures, this KPI quantifies how effectively the
schedule exploits available structural mitigation.

\medskip
\kpilabel{Formula.}
\begin{align*}
  \mathrm{CCR} &= \frac{\bigl|\{m \in M :
      \mathrm{WBGT}_m \geq 28\,^\circ\text{C}
      \;\text{and}\; v(m) \in R\}\bigr|}
    {\bigl|\{m \in M : \mathrm{WBGT}_m \geq
      28\,^\circ\text{C}\}\bigr|}
  \;\in\; [0,1].
\end{align*}
Complementary unmitigated exposure:
\begin{align*}
  \mathrm{UHE} &= \sum_{\substack{m \in M \\ v(m)\notin R}}
    \max\!\bigl(0,\;\mathrm{WBGT}_m - 28\bigr).
\end{align*}
\vspace{-8pt}

\noindent
Objective: $\;\displaystyle\max_{S}\;\mathrm{CCR}$.

\medskip
\kpilabel{Reference.}
\citet{FIFPRO2025}; \citet{Liljegren2008};
\citet{WorldWeatherAttribution2026}.
CCR definition is original (this work).
\end{kpibox}

% ----- KPI 2.4 -----
\begin{kpibox}{KPI~2.4 \textendash{} Altitude Disruption Index}

\kpilabel{Description.}
Under the base-camp model, each match $k$ requires a round trip
$b_i \!\to\! v_{i,k} \!\to\! b_i$.  The acclimatisation demand is
therefore proportional to the elevation differential between the base
camp and the match venue.  Mexico City ($\approx$2{,}240~m) and
Guadalajara ($\approx$1{,}560~m) versus sea-level US venues impose
an acclimatisation burden not captured by distance or rest-day metrics.
A threshold of 500~m is applied before the penalty begins, reflecting
the negligible physiological impact of modest elevation differences;
above this, the penalty increases linearly.  The sum runs from $k=1$
to $k=3$ since every match generates an independent round trip under
the base-camp model, analogously to KPI~1.1.

\medskip
\kpilabel{Formula.}
Define the penalisation function
\begin{align*}
  f(\Delta e) &= \frac{\max(0,\;\Delta e - 500)}{1000},
  \qquad \Delta e \text{ in metres,}
\end{align*}
so the penalty is zero below 500~m and increases linearly above it.
Then:
\begin{align*}
  \mathrm{ALT}_i &= \sum_{k=1}^{3}
    f\!\bigl(|\mathrm{elev}(v_{i,k}) - \mathrm{elev}(b_i)|\bigr).
\end{align*}
\vspace{-8pt}

\noindent
When $\mathrm{elev}(b_i)$ is unavailable, the median elevation of
FIFA-approved base-camp cities for team $i$'s association may be used
as a conservative estimate.
Objective: $\;\displaystyle\min_{S}\;\max_{i \in T}\,\mathrm{ALT}_i$.

\medskip
\kpilabel{Reference.}
FIFA Sports Medicine Commission (altitude guidelines);
\citet{MullanEtAl2025}.
Penalisation function $f(\cdot)$ and base-camp adaptation are original
(this work).
\end{kpibox}

% =================================================================
%  5.  CRITERION 3: COMPETITIVE INTEGRITY
% =================================================================
\section{Criterion 3: Competitive Integrity}

\subsection{Background: Collusion Risk in Group Tournaments}
\label{sec:collusion}

\textbf{What is collusion?}
In a round-robin group stage, the teams that play the \emph{final} match
of their group already know the complete standings table before kick-off.
If a particular result---one that may differ from the honest-play
outcome---allows \emph{both} teams to advance simultaneously, there is
an incentive to engineer that result cooperatively.  The most notorious
historical instance is the 1982 FIFA World Cup ``Disgrace of
Gij\'{o}n'': in the final round of Group~2, West Germany and Austria
played a second match knowing that a 1--0 German win would eliminate
Algeria despite Algeria having beaten West Germany.  The match ended
with exactly that scoreline after minimal on-field effort from both
sides.

\textbf{Why the $16 \times 3$ format was scrapped.}
FIFA's original plan for the 48-team tournament used 16 groups of three
teams.  \citet{Guyon2020} showed rigorously that this format generates
a collusion probability far exceeding that of any prior World Cup
because every group's final match is a two-team game where both clubs
know the standings after only two played matches.  Following widespread
academic and public criticism, FIFA in March 2023 adopted the
$12 \times 4$ format~\citep{FIFA2023Format}.
\citet{GuajardoKrumer2023, GuajardoKrumer2024} propose concrete schedule
constructions for this format that satisfy non-collusion, rest symmetry,
and no dead-rubbers simultaneously.

\textbf{The structural fix.}
Playing both matches of a group's final round \emph{simultaneously}
means no team can condition its behaviour on the real-time result of
the other match.  This is a hard scheduling constraint, not merely a
soft objective.

\subsection{Background: Dead Rubbers and Stakeless Matches}
\label{sec:deadrubber}

A match is \emph{stakeless} (colloquially a ``dead rubber'') if at
least one team is indifferent to the result---because it has already
secured qualification, is already eliminated, or cannot alter its
group-stage position regardless of the outcome.  Stakeless matches
undermine sporting integrity because teams lacking incentives may field
weakened line-ups or exert reduced effort, producing results that
distort the third-place rankings used to select additional qualifiers.

\citet{CsatoEtAl2024} and \citet{CsatoGyimesi2025} distinguish three
severity levels:
\begin{itemize}[noitemsep, topsep=4pt]
  \item \emph{Competitive}: both teams have a meaningful stake.
  \item \emph{Weakly stakeless}: one team is indifferent.
  \item \emph{Strongly stakeless}: both teams are indifferent.
\end{itemize}
These are weighted $0$, $0.5$, and $1.0$ respectively in KPI~3.2.

\subsection{Background: First-Mover Advantage in Round-Robins}
\label{sec:firstmover}

In a four-team round-robin group, the schedule specifies not only
\emph{which} teams play each other, but the \emph{order} in which they
do so across three rounds.  Each round consists of two matches; within
each round, one match is played first (earlier kickoff) and one second
(later kickoff).  \citet{KrumerEtAl2017} prove, in a three-player
sequential round-robin model where each match is a complete-information
all-pay auction, that the team competing in the \emph{first match of
the first round} has the highest subgame-perfect equilibrium win
probability (0.35 vs.\ 0.30 for the player who does not participate in
Round~1 at all).  The mechanism is informational: a strong early result
reduces uncertainty and shifts the strategic incentives of later rounds
in favour of the team that opened well.

For the 2026 four-team group stage, the final round is played
simultaneously (KPI~3.1), eliminating the first-mover effect in
Round~3.  The ordering advantage is therefore confined to Rounds~1
and~2, where one match within each round precedes the other and the
earlier-playing team observes fewer pending results when deciding its
strategic effort level.

\subsection{KPIs}

% ----- KPI 3.1 -----
\begin{kpibox}{KPI~3.1 \textendash{} Final-Round Simultaneity Compliance}

\kpilabel{Description.}
A binary indicator confirming whether each group's two final-round
matches kick off at the same time, the structural safeguard against
collusion established in the $12 \times 4$ format
(Section~\ref{sec:collusion}).  This KPI is primarily a compliance
verification measure; it should take its maximum value of~12 under any
valid 2026 schedule, and any deviation constitutes a critical design
flaw.

\medskip
\kpilabel{Formula.}
For group $g$ with final-round matches $m_1$ and $m_2$:
\begin{align*}
  \mathrm{SIM}_g &=
    \mathbf{1}\!\bigl[s(m_1) = s(m_2)\bigr].
\end{align*}
Tournament compliance score:
$\;\mathrm{SIM}_{\mathrm{tot}} =
\displaystyle\sum_{g \in G}\mathrm{SIM}_g \;\in\;\{0,\ldots,12\}$.
\\[4pt]
Hard constraint: $\mathrm{SIM}_g = 1$ for all $g \in G$.

\medskip
\kpilabel{Reference.}
\citet{Guyon2020}; \citet{GuajardoKrumer2024}; \citet{FIFA2023Format}.
\end{kpibox}

% ----- KPI 3.2 -----
\begin{kpibox}{KPI~3.2 \textendash{} Expected Stakeless Match Rate}

\kpilabel{Description.}
The probability-weighted fraction of final-round matches that are
stakeless (Section~\ref{sec:deadrubber}), computed by Monte Carlo
simulation of group outcomes using FIFA/Elo team-strength ratings.
The three-class weighting of \citet{CsatoGyimesi2025} penalises strongly
stakeless matches twice as heavily as weakly stakeless ones.

\medskip
\kpilabel{Formula.}
Simulate $N_{\mathrm{sim}}$ group-stage outcomes from an Elo-based
match-result model.  For each final-round match $m$ and each
replicate $r$, classify the match as competitive ($\mathcal{C}$),
weakly stakeless ($\mathcal{W}$), or strongly stakeless ($\mathcal{S}$).
The weighted stakeless rate is:
\begin{align*}
  \mathrm{WSL} &= \frac{1}{N_{\mathrm{sim}}}
    \sum_{r=1}^{N_{\mathrm{sim}}}
    \frac{\displaystyle\sum_{m \in M_{\mathrm{final}}}
      \Bigl[
        0 \cdot \mathbf{1}_{\mathcal{C}}^{(r)}(m)
        + 0.5 \cdot \mathbf{1}_{\mathcal{W}}^{(r)}(m)
        + 1 \cdot \mathbf{1}_{\mathcal{S}}^{(r)}(m)
      \Bigr]}
    {|M_{\mathrm{final}}|},
\end{align*}
where $M_{\mathrm{final}}$ is the set of third-round group matches
($|M_{\mathrm{final}}| = 24$).\\[4pt]
Objective: $\;\displaystyle\min_{S}\;\mathrm{WSL}$.

\medskip
\kpilabel{Reference.}
\citet{CsatoEtAl2024}; \citet{CsatoGyimesi2025} (weighting scheme
$0/0.5/1$); \citet{GuajardoKrumer2024}.
\end{kpibox}

% ----- KPI 3.3 -----
\begin{kpibox}{KPI~3.3 \textendash{} Round-Order Balance Index (First-Mover)}

\kpilabel{Description.}
Within a four-team round-robin group, the ordering of matches within
each round confers a systematic advantage on teams that play the earlier
kickoff (Section~\ref{sec:firstmover}).  \citet{KrumerEtAl2017} prove,
for the three-player case, that the first-match competitor has the
highest subgame-perfect equilibrium win probability.  For the 2026
four-team group, only Rounds~1 and~2 are sequentially ordered; Round~3
is simultaneous (KPI~3.1) and generates no first-mover effect.  This
KPI measures the dispersion of round-order advantage scores across all
48 teams: a lower value indicates that no subset of teams is
systematically favoured by the kickoff ordering.

\medskip
\kpilabel{Formula.}
For each group $g$ and each non-simultaneous round $r \in \{1,2\}$,
let $m_{g,r}^{(1)}$ and $m_{g,r}^{(2)}$ denote the two matches in
that round ordered by kickoff time:
$s\!\bigl(m_{g,r}^{(1)}\bigr) < s\!\bigl(m_{g,r}^{(2)}\bigr)$.
Define the binary first-mover indicator for team $i \in g$:
\begin{align*}
  \mathrm{pos}(i,r) &= \mathbf{1}\!\Bigl[
    \text{team } i \text{ plays in }
    m_{g(i),\,r}^{(1)}\Bigr]
  \;\in\; \{0,1\},
\end{align*}
so $\mathrm{pos}(i,r)=1$ if team $i$ plays the \emph{earlier} match
of round $r$ and $\mathrm{pos}(i,r)=0$ otherwise.  The positional
advantage score is:
\begin{align*}
  a_i &= \sum_{r=1}^{2} \mathrm{pos}(i,r) \;\in\; \{0,1,2\}.
\end{align*}
$a_i = 2$: team plays first in both non-simultaneous rounds (maximum
first-mover benefit).  $a_i = 0$: team plays second in both (minimal
benefit).  A perfectly balanced schedule gives $a_i = 1$ for all
$i \in T$, yielding:
\begin{align*}
  \mathrm{FMB} &= \sigma\!\bigl(\{a_i : i \in T\}\bigr) = 0.
\end{align*}
Objective: $\;\displaystyle\min_{S}\;\mathrm{FMB}$.

\medskip
\kpilabel{Reference.}
\citet{KrumerEtAl2017}; \citet{CsatoEtAl2024}.  See also
\citet{Russell1980} (carry-over effect value) and
\citet{CavdarogluAtan2022} (integrated scheduling model).
Definition of $\mathrm{pos}(i,r)$ and extension to four-team groups
are original (this work).
\end{kpibox}

% =================================================================
%  6.  CRITERION 4: OPERATIONAL EFFICIENCY
% =================================================================
\section{Criterion 4: Operational Efficiency}

\subsection{Overview}

Operational efficiency concerns the logistical load on the 16 host
venues and the convenience of in-person attendance.  Venue
overload---too many matches within a short window---degrades pitch
quality, exhausts grounds-crew and security capacity, and strains local
transport infrastructure.  Same-city scheduling conflicts split local
resources and create confusion for fans.  These concerns are secondary
to welfare and integrity, but they affect the experience of the
approximately 6.5~million expected group-stage attendees and the host
cities' ability to fulfil their delivery commitments.

\subsection{KPIs}

% ----- KPI 4.1 -----
\begin{kpibox}{KPI~4.1 \textendash{} Venue-Load Balance}

\kpilabel{Description.}
The evenness of the distribution of group-stage matches across the 16
host venues.  A schedule that concentrates fixtures in a few stadiums
while leaving others underused creates operational strain at the busy
venues and reduces economic impact at the quiet ones.  The coefficient
of variation (CV) provides a scale-free measure of dispersion.

\medskip
\kpilabel{Formula.}
Let $n_v$ be the number of group-stage matches at venue $v$.  Then:
\begin{align*}
  \mathrm{VLB} &= \frac{\sigma(\{n_v : v \in V\})}
                       {\mu(\{n_v : v \in V\})}.
\end{align*}
\vspace{-8pt}

\noindent
Objective: $\;\displaystyle\min_{S}\;\mathrm{VLB}$.

\medskip
\kpilabel{Reference.}
\citet{RibeiroUrrutia2007} (venue-load balancing objectives in
tournament scheduling).  Application to the 16-venue 2026 format is
original (this work).
\end{kpibox}

% ----- KPI 4.2 -----
\begin{kpibox}{KPI~4.2 \textendash{} Fan Accessibility and Same-City Overlap}

\kpilabel{Description.}
Captures two related operational concerns.  First,
\emph{same-city overlap} counts instances where two matches occur in
the same city on the same day, potentially straining transport, security,
and hospitality resources beyond single-match capacity.  Second, the
\emph{fan accessibility index} measures the feasibility of supporters
following a team to all three group-stage venues within a geographically
reasonable distance threshold.

\medskip
\kpilabel{Formula.}
Same-city overlap count:
\begin{align*}
  \mathrm{SCO} &= \sum_{\substack{m,m'\in M \\ m < m'}}
    \mathbf{1}\!\bigl[\mathrm{city}(m)=\mathrm{city}(m')
    \;\wedge\; d(m)=d(m')\bigr].
\end{align*}
Fan accessibility index for team $i$, with threshold $\bar{d}$ (e.g.,
$\bar{d}=500$~km, approximately one day's drive):
\begin{align*}
  \mathrm{FA}_i &= \sum_{k=1}^{2}
    \mathbf{1}\!\bigl[\mathrm{dist}(v_{i,k},v_{i,k+1})
    \leq \bar{d}\bigr].
\end{align*}
\vspace{-8pt}

\noindent
Objective: $\;\displaystyle\min_{S}\;\mathrm{SCO}$.

\medskip
\kpilabel{Reference.}
\citet{KendallEtAl2010}; \citet{RibeiroUrrutia2007}.
Application is original (this work).
\end{kpibox}

% =================================================================
%  7.  CRITERION 5: BROADCAST ATTRACTIVENESS
% =================================================================
\section{Criterion 5: Broadcast Attractiveness and Commercial Value}

\subsection{Overview}

FIFA's broadcast rights for the 2026 World Cup are projected to exceed
\$11~billion in total revenue.  Kickoff-time allocation is consequently
one of the most commercially sensitive scheduling decisions.  The primary
driver is European prime time (19:00--23:00 CET), which simultaneously
represents the world's largest football audience and---for US East Coast
venues---a midday kickoff that coincides with peak heat.  This criterion
explicitly models broadcast and commercial value as potential antagonists
of Criteria~1 and~2, enabling a transparent Pareto trade-off analysis.

\subsection{KPIs}

% ----- KPI 5.1 -----
\begin{kpibox}{KPI~5.1 \textendash{} Prime-Time Alignment Score}

\kpilabel{Description.}
The audience-weighted share of group-stage match hours that fall within
the prime-time window of key broadcasting markets.
\citet{RibeiroUrrutia2007} formalise broadcast objectives jointly with
scheduling fairness in the context of the Brazilian national league,
providing the methodological foundation for this KPI.

\medskip
\kpilabel{Formula.}
Let $\mathcal{R}$ be the set of key markets, $\mathrm{pop}_r$ the
audience weight of market $r$, and $w_r(s) \in [0,1]$ the indicator
(or fractional overlap) that kickoff slot $s$ falls within the
prime-time window of market $r$:
\begin{align*}
  \mathrm{PT} &= \sum_{m \in M}\;\sum_{r \in \mathcal{R}}\;
    \mathrm{pop}_r \cdot w_r\!\bigl(s(m)\bigr).
\end{align*}
\vspace{-8pt}

\noindent
Objective: $\;\displaystyle\max_{S}\;\mathrm{PT}$,
\emph{subject to} heat and rest constraints from C1 and C2.

\medskip
\kpilabel{Reference.}
\citet{RibeiroUrrutia2007}.
\end{kpibox}

% ----- KPI 5.2 -----
\begin{kpibox}{KPI~5.2 \textendash{} Marquee-Match Slot Quality and Overlap Penalty}

\kpilabel{Description.}
Not all matches carry equal broadcast value.  Matches between high-ranked
teams, group deciders in the final round, and high-interest regional
rivalries command disproportionate audience attention.  This KPI
combines (a)~the global slot quality of high-profile matches and
(b)~a penalty for scheduling two high-profile matches simultaneously,
which splits the global audience~\citep{RibeiroUrrutia2007}.

\medskip
\kpilabel{Formula.}
Define the slot quality $q(s) =
\sum_{r \in \mathcal{R}} \mathrm{pop}_r \cdot w_r(s)$ (as in KPI~5.1).
For a match between teams $i$ and $j$, define a match-value score
$\mu_{ij} = (E_i^{-1} + E_j^{-1})/2$ (higher for stronger teams).
Then the marquee-match slot quality is:
\begin{align*}
  \mathrm{MSQ} &= \sum_{m \in M} \mu_{ij(m)} \cdot
    q\!\bigl(s(m)\bigr).
\end{align*}
Simultaneous high-profile match penalty ($P \subseteq M$):
\begin{align*}
  \mathrm{OVL} &= \sum_{\substack{m,m'\in P \\ m < m'}}
    \mathbf{1}\!\bigl[s(m) = s(m')\bigr].
\end{align*}
Net objective: $\;\displaystyle\max_{S}\;
\bigl(\mathrm{MSQ} - \lambda \cdot \mathrm{OVL}\bigr)$,
\quad $\lambda > 0$.

\medskip
\kpilabel{Reference.}
\citet{RibeiroUrrutia2007}; \citet{GuajardoKrumer2024}.
Original composite (this work).
\end{kpibox}

% ----- KPI 5.3 -----
\begin{kpibox}{KPI~5.3 \textendash{} Host-City Economic Equity}

\kpilabel{Description.}
Host cities bear significant upfront costs---security, infrastructure,
permitting---while FIFA retains broadcast and ticketing revenue.  Their
economic return depends heavily on the commercial quality (prime-time
status, high-profile matches) of fixtures allocated to them.  This KPI
measures how equitably high-value matches are distributed across the 16
venues using the Gini coefficient as a normalised dispersion measure.

\medskip
\kpilabel{Formula.}
Define the commercial value of venue $v$:
$\mathrm{VC}_v = \sum_{m:\,v(m)=v}
q\!\bigl(s(m)\bigr)\cdot\mu_{ij(m)}$.  Then:
\begin{align*}
  \mathrm{HCE} &= 1 \;-\;
    \mathrm{Gini}\!\bigl(\{\mathrm{VC}_v : v \in V\}\bigr)
  \;\in\; [0,1].
\end{align*}
\vspace{-8pt}

\noindent
A value of~1 indicates perfectly equal commercial value distribution.
Objective: $\;\displaystyle\max_{S}\;\mathrm{HCE}$.

\medskip
\kpilabel{Reference.}
\citet{RibeiroUrrutia2007} (fairness of match allocation).
Gini coefficient: standard reference.  Application is original
(this work).
\end{kpibox}

% =================================================================
%  8.  COMPOSITE / CROSS-CRITERIA KPIs
% =================================================================
\section{Composite and Cross-Criteria KPIs}
\label{sec:composite}

Single-criterion KPIs support granular analysis but can obscure the
overall burden borne by a team or the global fairness of the schedule.
The five composite KPIs below aggregate across criteria to provide
integrative measures.  All normalise
component KPIs to $[0,1]$ to enable principled trade-off weighting.

% ----- C1 -----
\begin{kpibox}{C1 \textendash{} Recovery-Adjusted Hardship Index (RAHI)}

\kpilabel{Description.}
A single per-team score combining the four binding welfare burdens:
travel distance (KPI~1.1, formulated under the base-camp round-trip
model), jet-lag cost (KPI~1.3), heat load (KPI~2.2), and rest
(KPI~1.5, entering negatively because more rest reduces hardship).
All components are normalised to $[0,1]$ and combined via
user-specified weights.  This is the flagship single-objective measure
for comparing alternative schedules.

\medskip
\kpilabel{Formula.}
\begin{align*}
  \mathrm{RAHI}_i \;=\;
    \alpha \cdot \frac{\mathrm{TD}_i}{\max_{i'}\mathrm{TD}_{i'}}
    + \beta \cdot \frac{\mathrm{JL}_i}{\max_{i'}\mathrm{JL}_{i'}}
    + \gamma \cdot \frac{\mathrm{HL}_i}{\max_{i'}\mathrm{HL}_{i'}}
    - \delta \cdot \frac{\sum_k g_{i,k}}
                        {\max_{i'}\sum_k g_{i',k}},
\end{align*}
with $\alpha,\beta,\gamma,\delta \geq 0$ and
$\alpha + \beta + \gamma + \delta = 1$.
Default weights: $\alpha = 0.30$, $\beta = 0.20$, $\gamma = 0.30$,
$\delta = 0.20$.\\[4pt]
Minimax objective (protect the worst-off team):
$\;\displaystyle\min_{S}\;\max_{i \in T}\;\mathrm{RAHI}_i$.

\medskip
\kpilabel{Reference.}
Components: \citet{Easton2001} (TD); \citet{Recht1995} (JL);
\citet{Liljegren2008}, \citet{FIFPRO2025} (HL);
\citet{AtanCavdaroglu2018} (rest).
Composite formulation is original (this work).
\end{kpibox}

% ----- C2 -----
\begin{kpibox}{C2 \textendash{} Schedule Fairness Index (SFI)}

\kpilabel{Description.}
Converts the per-team RAHI into a global distributional fairness
statistic.  Jain's fairness index~\citep{Jain1984} was originally
developed for network resource allocation and has since been widely
adopted in computational fairness literature.  A value of~1 indicates
all teams bear identical hardship; values below~1 signal concentration
of hardship in a subset of teams.

\medskip
\kpilabel{Formula.}
\begin{align*}
  \mathrm{SFI} \;=\;
    \frac{\Bigl(\displaystyle\sum_{i \in T}
      \mathrm{RAHI}_i\Bigr)^2}
         {|T|\cdot\displaystyle\sum_{i \in T} \mathrm{RAHI}_i^2}
  \;\in\;\Bigl[\tfrac{1}{|T|},\;1\Bigr].
\end{align*}
\vspace{-8pt}

\noindent
Objective: $\;\displaystyle\max_{S}\;\mathrm{SFI}$.

\medskip
\kpilabel{Reference.}
\citet{Jain1984}; \citet{GuajardoEtAl2026} (travel fairness).
Application to schedule hardship is original (this work).
\end{kpibox}

% ----- C3 -----
\begin{kpibox}{C3 \textendash{} Heat-Roof Mitigation Gap (HRMG)}

\kpilabel{Description.}
Quantifies avoidable heat exposure: the additional WBGT stress that
could be eliminated by reassigning safety-threshold-exceeding matches to
roofed venues or cooler (evening) slots, given the available covered
infrastructure.  A large HRMG indicates the current schedule is failing
to exploit structural mitigation, providing direct actionable insight
for the optimised schedule.

\medskip
\kpilabel{Formula.}
Let $\mathrm{WBGT}^*_m$ denote the counterfactual WBGT if match $m$
were assigned to the best available roofed venue (using
equation~\eqref{eq:wbgt_indoor}) or the latest evening slot on the same
date.  Then:
\begin{align*}
  \mathrm{HRMG} &= \sum_{\substack{m \in M \\
    \mathrm{WBGT}_m \geq 28}}
    \max\!\bigl(0,\;\mathrm{WBGT}_m - \mathrm{WBGT}^*_m\bigr).
\end{align*}
\vspace{-8pt}

\noindent
Objective: $\;\displaystyle\min_{S}\;\mathrm{HRMG}$.

\medskip
\kpilabel{Reference.}
\citet{FIFPRO2025}; \citet{Liljegren2008};
\citet{WorldWeatherAttribution2026}.
Composite formulation is original (this work).
\end{kpibox}

% ----- C4 -----
\begin{kpibox}{C4 \textendash{} Welfare-Broadcast Tension Index (WBT)}

\kpilabel{Description.}
Explicitly quantifies the conflict between broadcast value (C5) and
player welfare (C2): how much of the prime-time commercial value is
generated by matches that breach the $28\,^\circ$C safety threshold.
A schedule achieving high broadcast value while keeping matches in safe
windows scores low; one that sacrifices welfare for TV ratings scores
high.  This index supports a transparent Pareto analysis.

\medskip
\kpilabel{Formula.}
\begin{align*}
  \mathrm{WBT} &= \sum_{m \in M}
    q\!\bigl(s(m)\bigr)\cdot
    \mathbf{1}\!\bigl[\mathrm{WBGT}_m \geq 28\,^\circ\text{C}\bigr].
\end{align*}
\vspace{-8pt}

\noindent
Objective: $\;\displaystyle\min_{S}\;\mathrm{WBT}$,
subject to an acceptable lower bound on $\mathrm{PT}$ from KPI~5.1.

\medskip
\kpilabel{Reference.}
\citet{RibeiroUrrutia2007} (broadcast--fairness bicriteria framing);
\citet{FIFPRO2025}.  Composite formulation is original (this work).
\end{kpibox}

% ----- C5 -----
\begin{kpibox}{C5 \textendash{} Disruption-Robust Worst-Case Hardship (DRWH)}

\kpilabel{Description.}
The worst-team RAHI under the most adverse feasible single-match one-day
perturbation.  North American summer conditions---severe thunderstorms,
heat-related postponements---can force one-day slippage of individual
fixtures.  This KPI asks: if any single match is displaced by one day,
how much does the hardship of the worst-off team increase?  A small
robustness premium indicates a schedule that degrades gracefully under
disruption.

\medskip
\kpilabel{Formula.}
\begin{align*}
  \mathrm{DRWH} &= \max_{i \in T}\;
    \max_{\delta \in \{-1,+1\},\;m \in M}\;
    \mathrm{RAHI}_i\!\big|_{d(m)\,\leftarrow\,d(m)+\delta}.
\end{align*}
Robustness premium:
$\;\mathrm{RP} = \mathrm{DRWH} - \max_{i \in T}\,\mathrm{RAHI}_i$.
\\[4pt]
Objective: $\;\displaystyle\min_{S}\;\mathrm{RP}$.

\medskip
\kpilabel{Reference.}
Recoverable robustness: \citet{Liebchen2009};
\citet{KendallEtAl2010}.  Application is original (this work).
\end{kpibox}

% =================================================================
%  9.  ROBUSTNESS TO DISRUPTIONS
% =================================================================
\section{Robustness to Disruptions}

Disruption risks in the 2026 group stage include severe thunderstorms
(especially in the Central region), heat-related postponements at
open-air venues, and potential flight delays from security or weather
events.  A robust schedule is one in which a single-match one-day slip
cannot violate any hard welfare constraint, particularly the rest floor.
The single KPI below serves as a hygiene check: it should be satisfied
by any feasible schedule but is worth verifying explicitly given the
dense 17-day window.

\begin{kpibox}{R1 \textendash{} Disruption-Robust Rest Margin}

\kpilabel{Description.}
The minimum rest available to any team under the worst feasible
single-match one-day disruption.  If this value remains at or above the
hard minimum of $g_{\min}=2$ rest days, the schedule is
\emph{rest-robust}.

\medskip
\kpilabel{Formula.}
\begin{align*}
  \mathrm{RM} &= \min_{i \in T}\;\min_{k \in \{1,2\}}\;
    \min_{\delta \in \{-1,+1\},\;m \in M}\;
    g_{i,k}\!\big|_{d(m)\,\leftarrow\,d(m)+\delta}.
\end{align*}
\vspace{-8pt}

\noindent
Feasibility condition: $\;\mathrm{RM} \geq g_{\min} = 2$.

\medskip
\kpilabel{Reference.}
\citet{Liebchen2009} (recoverable robustness).
Application is original (this work).
\end{kpibox}

% =================================================================
%  10.  SUMMARY TABLE
% =================================================================
\section{Summary of the KPI Framework}
\label{sec:summary}

Table~\ref{tab:summary} provides a consolidated reference for all
25~KPIs defined in this report.

\setlength{\LTpre}{8pt}
\setlength{\LTpost}{6pt}

\begin{longtable}{@{}p{0.9cm} p{4.6cm} p{2.8cm} p{1.2cm} p{4.0cm}@{}}
\caption{Complete KPI framework for the 2026 FIFA World Cup
         group-stage schedule evaluation.}
\label{tab:summary}\\
\toprule
\textbf{ID} & \textbf{Name} & \textbf{Criterion} &
\textbf{Dir.} & \textbf{Primary Reference} \\
\midrule
\endfirsthead
\multicolumn{5}{c}{\small\itshape Table~\ref{tab:summary}
  \textendash{} continued from previous page.} \\[4pt]
\toprule
\textbf{ID} & \textbf{Name} & \textbf{Criterion} &
\textbf{Dir.} & \textbf{Primary Reference} \\
\midrule
\endhead
\midrule
\multicolumn{5}{r}{\small\itshape Continued on next page.} \\
\endfoot
\bottomrule
\endlastfoot

1.1 & Total Travel Distance & Travel/Rest & $\min$ &
      \citet{Easton2001} \\[2pt]
1.2 & Intra-Group Travel Dispersion & Travel/Rest & $\min$ &
      \citet{GuajardoEtAl2026} \\[2pt]
1.3 & Circadian Shift Cost & Travel/Rest & $\min$ &
      \citet{Recht1995} \\[2pt]
1.4 & Match-Venue Geographic Dispersion & Travel/Rest & $\min$ &
      \citet{MikamiTravel2026} \\[2pt]
1.5 & Per-Team Rest Adequacy & Travel/Rest & $\max$ &
      \citet{AtanCavdaroglu2018} \\[2pt]
1.6 & Rest Asymmetry Between Opponents & Travel/Rest & $\min$ &
      \citet{AtanCavdaroglu2018} \\[2pt]
1.7 & Entry \& Visa Restriction Exposure & Travel/Rest & $\min$ &
      Original (this work) \\
\midrule
2.1 & Tournament WBGT Score & Heat & $\min$ &
      \citet{Liljegren2008} \\[2pt]
2.2 & Per-Team Heat Load & Heat & $\min$ &
      \citet{FIFPRO2025} \\[2pt]
2.3 & Roof Conditioning Ratio & Heat & $\max$ &
      \citet{FIFPRO2025} \\[2pt]
2.4 & Altitude Disruption Index & Heat & $\min$ &
      Original (this work) \\
\midrule
3.1 & Final-Round Simultaneity & Comp.\ Integrity & $\max$ &
      \citet{Guyon2020} \\[2pt]
3.2 & Expected Stakeless Match Rate & Comp.\ Integrity & $\min$ &
      \citet{CsatoEtAl2024} \\[2pt]
3.3 & Round-Order Balance Index & Comp.\ Integrity & $\min$ &
      \citet{KrumerEtAl2017} \\
\midrule
4.1 & Venue-Load Balance & Operational & $\min$ &
      \citet{RibeiroUrrutia2007} \\[2pt]
4.2 & Fan Accessibility \& Same-City Overlap & Operational & $\min$ &
      Original (this work) \\
\midrule
5.1 & Prime-Time Alignment Score & Broadcast & $\max$ &
      \citet{RibeiroUrrutia2007} \\[2pt]
5.2 & Marquee-Match Slot Quality & Broadcast & $\max$ &
      \citet{RibeiroUrrutia2007} \\[2pt]
5.3 & Host-City Economic Equity & Broadcast & $\max$ &
      Original (this work) \\
\midrule
C1  & Recovery-Adjusted Hardship Index & Composite & $\min$ &
      Original (this work) \\[2pt]
C2  & Schedule Fairness Index & Composite & $\max$ &
      \citet{Jain1984} \\[2pt]
C3  & Heat-Roof Mitigation Gap & Composite & $\min$ &
      Original (this work) \\[2pt]
C4  & Welfare-Broadcast Tension Index & Composite & $\min$ &
      Original (this work) \\[2pt]
C5  & Disruption-Robust Worst-Case Hardship & Composite & $\min$ &
      Original (this work) \\
\midrule
R1  & Disruption-Robust Rest Margin & Robustness & feasib. &
      \citet{Liebchen2009} \\
\end{longtable}

% =================================================================
%  BIBLIOGRAPHY
% =================================================================
\bibliographystyle{abbrvnat}
\bibliography{references}

\end{document}